\scanpage{8}
a aussi l'inégalité suivante (résultant déjà de la suite exacte des trois derniers termes de (8.5.1)):
\[
\dim(E(Y)\otimes_W k)_{ss}\handwritten{\geq}\dim E(Y,\underline O_Y)_{ss}
\deleted{-\dim\mathrm{Ker}\,w}
\tag*{(8.5.6)}
\]
\[
\deleted{\dim E(Y,\underline O_Y)_{ss}-t^{n-1}(X,p)\;,}
\]
\deleted{où $t^i(X,p)$ désigne le nombre de \inserted{$p$-}coéfficients de torsion de $X$ (défini via la cohomologie cristalline de $X$ en dimension $i$).}

\uline{Remarques} 8.5.7. Je pense que l'homomorphisme $v$ de (8.5.2) est injectif\deleted{\unsure{,}} (ce qui simplifie en conséquence (8.5.1) et (8.5.5), puisque on aurait $\mathrm{Ker}\,v=0$ ); c'est vrai en tous cas si \inserted{\handwritten{$n=2$, ou si}} $Y$ est \deleted{de degré \unsure{assez|asse}} u\deleted{\unsure{j}}ne section par une hypersurface de degré $d$ assez grand pour qu'on ait
\[
H^{n-1}(X,\underline O_X(-d))=0\;.
\]

\deleted{D'autre part, on peut si on veut dans 8.5.5 \deleted{\unsure{bolquer}}bloquer ensemble les deux termes $\dim\mathrm{Coker}\,w-\dim\mathrm{Ker}\,w$, \deleted{car on aura comme} \inserted{pour trouver} la différence de rangs \deleted{$\dim\mathrm{Coker}\,w-\dim\mathrm{Ker}\,w$ \unsure{est}} des deux membres de (8.5.3).}

8.5.8. Lorsque \deleted{dans} l'on applique (8.5.5) en prenant pour $Y$ une section de $X$ avec une hypersurface de degré $d$ ``assez générale'' (de telle façon que le rang semi-simple de $E(Y,\underline O_Y)$ ait la valeur maxima), alors on trouve qu'au\deleted{\unsure{t}} terme $-\dim\mathrm{Ker}\,v$ près (sans doute toujours nul, cf. 8.5.\deleted{\unsure{6}}7) le deuxième membre est \uline{indépendant du choix} \deleted{de} \uline{particulier de \handwritten{$Y$}}\unsure{.} Cette remarque s'applique en particulier (8.5.7) si $d$ est assez grand. On en conclut que si $k$ est fini, \deleted{\unsure{les,}} on peut trouver un ouvert \inserted{$U_d$} de l'espace projectif des hypersurfaces de degré $d$, tel que pour \deleted{tout} \inserted{un} point \inserted{variable} de $U$ à coéfficients dans une extension finie $k'$ de $k$, le nombre de valeurs propres de Frobenius opérant sur la cohomologie évanescente de $Y_{\overline{k'}}$ et qui sont des unités $p$-adiques soit constant. \deleted{Mieux, quitte à agrandir un peut $k$, de sorte que les valeurs propres de $\mathrm{Frob}_k$} Nous désirons mon%
% archive sheet number: 7
